\documentclass[11pt]{article}
\usepackage[margin=1in]{geometry}
\usepackage{amsmath,amssymb,amsthm}
\usepackage{booktabs}
\usepackage{multirow}
\usepackage{graphicx}
\usepackage{algorithm}
\usepackage{algpseudocode}
\newcommand{\bm}[1]{\boldsymbol{#1}}  % bm package absent in TinyTeX repo; \boldsymbol from amsmath
\usepackage[hidelinks]{hyperref}
\usepackage{xcolor}

\newcommand{\Phimat}{\bm{\Phi}}
\newcommand{\A}{\mathbf{A}}
\newcommand{\K}{\mathbf{K}}
\newcommand{\x}{\mathbf{x}}
\newcommand{\y}{\mathbf{y}}
\newcommand{\bmu}{\bm{\mu}}
\newcommand{\I}{\mathbf{I}}
\DeclareMathOperator*{\argmin}{arg\,min}

\title{\textbf{Scale-Matched Data Consistency for Non-Hot Diffusion Priors:\\
An HQS Framework for Linear Image Restoration}}
\date{\today}

\begin{document}
\maketitle

\begin{abstract}
We present an image-restoration method that couples a \emph{non-hot} diffusion prior
with a per-step \emph{Half-Quadratic Splitting} data-consistency
solver. The forward blur operator of the prior, $\K_t$, and the measurement operator $\A$
are simultaneously diagonalized in a fixed orthogonal transform $\Phimat$ (DCT for
Neumann/symmetric blur, DFT for periodic/asymmetric blur); this makes $\A\K_t=\K_t\A$ and
turns each HQS data step into a closed-form, per-frequency Wiener filter with \emph{no}
backpropagation through the network. We give the problem statement, the concrete HQS
derivation and update, the algorithm, and numerical results on held-out CelebA-HQ-256 for
Gaussian and motion deblurring, benchmarked against Diffusion Posterior Sampling (DPS).
\end{abstract}

\section{Problem definition}
\label{sec:problem}

\paragraph{Linear inverse problem.}
Let $\x_0\in\mathbb{R}^{d}$ ($d=3HW$) be a clean image. We observe
\begin{equation}
\y \;=\; \A\,\x_0 \;+\; \bm{n},\qquad \bm{n}\sim\mathcal{N}(\bm 0,\;\sigma_y^2\I),
\label{eq:model}
\end{equation}
where $\A$ is a known linear degradation and
$\sigma_y$ the noise level. The task is to recover $\x_0$ from $\y$.

\paragraph{Non-hot diffusion prior.}
A non-hot diffusion prior (cold diffusion / IHDM) defines a
\emph{deterministic} forward corruption by progressive blurring,
\begin{equation}
\x_t \;=\; \K_t\,\x_0,\qquad t=0,1,\dots,N,\qquad \K_0=\I,
\label{eq:forward}
\end{equation}
where $\K_t$ is a heat-dissipation (Gaussian-smoothing) operator of increasing scale
$\sigma_t$. A network $R_\theta$ is trained to invert one level of blur; its
reverse step produces a prior mean
\begin{equation}
\bmu_{t-1} \;=\; R_\theta(\x_t,\,t)\;\approx\;\mathbb{E}[\x_{t-1}\mid \x_t].
\label{eq:reverse}
\end{equation}
Crucially, $R_\theta$ operates entirely in the \emph{image domain}: it maps blurred pixels
to (a residual toward) less-blurred pixels; the transform basis enters only inside $\K_t$.

\paragraph{Shared eigenbasis.}
Both $\K_t$ and $\A$ are diagonalized by a fixed orthogonal transform $\Phimat$:
\begin{equation}
\K_t=\Phimat^{-1}\,\mathrm{diag}(\hat{k}_t)\,\Phimat,\qquad
\A=\Phimat^{-1}\,\mathrm{diag}(\hat{a})\,\Phimat,\qquad
\hat{k}_t[j]=e^{-\tfrac12\sigma_t^2\lambda_j},
\label{eq:diag}
\end{equation}
with $\lambda_j$ the Laplacian eigenvalues. The correct basis depends on boundary
condition and kernel symmetry:
\emph{(i)} \textbf{Neumann/symmetric} blur (Gaussian) $\Rightarrow$ $\Phimat=$ DCT,
$\hat a$ real; \emph{(ii)} \textbf{periodic/asymmetric} blur $\Rightarrow$
$\Phimat=$ DFT, $\hat a$ complex. When $\A$ and $\K_t$ share $\Phimat$
they commute,
\begin{equation}
\A\,\K_t \;=\; \K_t\,\A,
\label{eq:commute}
\end{equation}
which is the property our solver exploits.

\section{Method: per-step MAP via HQS}
\label{sec:method}

\paragraph{Restoration as guided reverse diffusion.}
We run the prior's reverse chain from $t=N$ to $0$ and, at every level, correct the prior
mean $\bmu_{t-1}$ toward the measurement. We interpret the correction as one step of a
\emph{plug-and-play} MAP estimate: with a Gaussian prior-proximal term of variance
$\delta^2$ around $\bmu_{t-1}$ and the Gaussian likelihood~\eqref{eq:model},
\begin{equation}
\x_{t-1} \;=\; \argmin_{\x}\;
\underbrace{\frac{1}{2\delta^2}\,\|\x-\bmu_{t-1}\|^2}_{\text{prior proximal}}
\;+\;
\underbrace{\frac{1}{2\sigma_y^2}\,\|\A\x-\y_{t-1}\|^2}_{\text{data consistency}} .
\label{eq:map}
\end{equation}
This is exactly the HQS subproblem obtained by variable-splitting the MAP objective
$\min_\x \tfrac{1}{2\sigma_y^2}\|\A\x-\y\|^2 + \mathcal{R}(\x)$ and replacing the prior
proximal operator by the learned diffusion reverse step. Substituting a \emph{classical}
denoiser for that prox---e.g.\ a total-variation (Chambolle) step---recovers standard
plug-and-play HQS; we report this as the \textbf{HQS$+$TV} baseline in Sec.~\ref{sec:results}.

\paragraph{Scale-matched measurement.}
Rather than driving every level with the raw $\y$, we \emph{scale-match} the measurement
to the current blur level. Because $\A$ and $\K_t$ commute~\eqref{eq:commute}, in the
noiseless case
\begin{equation}
\A\,\x_{t-1}\;=\;\A\,\K_{t-1}\,\x_0\;=\;\K_{t-1}\,\A\,\x_0\;=\;\K_{t-1}\,\y ,
\end{equation}
so the exact ``companion'' operator is $\mathbf{L}_{t-1}=\K_{t-1}$ and we set
\begin{equation}
\y_{t-1}\;=\;\mathbf{L}_{t-1}\,\y\;=\;\K_{t-1}\,\y .
\label{eq:scalematch}
\end{equation}
Early (large $t$) the target is heavily blurred, so data consistency is loose; as
$t\!\to\!0$, $\K_0=\I$ and $\y_0=\y$, enforcing the full measurement. This anneals data
consistency along the trajectory.

\paragraph{Closed form.}
Since the objective~\eqref{eq:map} is quadratic and \emph{both} terms are diagonal in
$\Phimat$,
the minimizer is available per frequency $j$ with no iteration:
\begin{equation}
\boxed{\;
\hat{\x}_{t-1}[j]
=\frac{w_p\,\hat{\bmu}_{t-1}[j]\;+\;w_y(t)\,\overline{\hat{a}[j]}\;\widehat{(\K_{t-1}\y)}[j]}
       {w_p\;+\;w_y(t)\,\lvert\hat{a}[j]\rvert^{2}} ,
\qquad \x_{t-1}=\Phimat^{-1}\hat{\x}_{t-1},
\;}
\label{eq:wiener}
\end{equation}
with weights
\begin{equation}
w_p=\frac{\alpha_p}{\delta^{2}},\qquad
w_y(t)=\frac{\alpha_y}{\sigma_y^{2}}\Bigl(1-\tfrac{t-1}{N}\Bigr),
\label{eq:weights}
\end{equation}
where $\overline{(\cdot)}$ is complex conjugation. The scalars $\alpha_p,\alpha_y>0$ are
dimensionless weights on the prior and data terms (\emph{not} the Laplacian eigenvalues
$\lambda_j$ of Eq.~\eqref{eq:diag}); only their ratio $\alpha_p/\alpha_y$ shapes the filter,
$\delta$ is the prior-proximal standard deviation, and $\sigma_y$ the noise level
(see the parameter sheet, Table~\ref{tab:params}).
Equation~\eqref{eq:wiener} is a per-frequency Wiener filter:
at frequencies with $|\hat a[j]|\!\approx\!0$ (e.g.\ spectral nulls of a motion kernel)
the estimate falls back to the prior $\hat{\bmu}[j]$, while at well-conditioned
frequencies it follows the data. The linear ramp in~\eqref{eq:weights} realizes the
``late'' schedule (data pull strongest near $t=0$). No gradient is taken through
$R_\theta$; the network is used only as the prior-proximal map producing $\bmu_{t-1}$.

\paragraph{Basis selection.}
For Gaussian/heat $\A$, $\Phimat=$ DCT and $\hat a$ is real. For motion $\A$, a DCT prior
does \emph{not} commute with $\A$ (complex phase cannot live in a real cosine basis); we
therefore evaluate the entire pipeline in the DFT basis, using the periodic heat symbol
$\hat k_t[j]=\exp(-\tfrac12\sigma_t^2\lvert\omega_j\rvert^2)$ with
$\omega_j=2\pi\,\mathrm{fftfreq}$, which restores $\A\K_t=\K_t\A$ and hence the exact
scale-match~\eqref{eq:scalematch} and closed form~\eqref{eq:wiener}.

\paragraph{Parameter settings.}
Table~\ref{tab:params} is the full parameter sheet. Only the ratio $\alpha_p/\alpha_y$
matters for the filter (effective regularization
$\tfrac{\alpha_p}{\alpha_y}\tfrac{\sigma_y^2}{\delta^2}$); we deliberately up-weight data
($\alpha_y\!\gg\!\alpha_p$) because the one-step diffusion reverse is only an approximate
proximal operator and the measurement must be pushed harder than a textbook MAP would set.

\begin{table}[h]
\centering
\caption{Parameter choice sheet. Top block: solver hyperparameters; bottom block: observation
model / evaluation.}
\label{tab:params}
\begin{tabular}{@{}ll p{6.6cm}@{}}
\toprule
Symbol & Meaning & Value \\
\midrule
$\alpha_p$   & prior-term weight        & $1.0$ \\
$\alpha_y$   & data-term weight         & $64$ \\
$\delta$     & prior-proximal std       & $0.01$ \\
$\sigma_y$   & measurement-noise std    & $0.05$ \\
$N$          & number of blur levels    & $200$ \\
$\A$         & degradation operator     & Gaussian $\sigma=4$ \,/\, Levin09~\#1 ($19\times19$) \\
noise        & additive Gaussian        & $\sigma_y=0.05$ \\
\bottomrule
\end{tabular}
\end{table}

\section{Algorithm}
\label{sec:algo}

\begin{algorithm}[h]
\caption{HQS restoration with a non-hot diffusion prior}
\label{alg:hqs}
\begin{algorithmic}[1]
\Require measurement $\y$; operator $\A$ (transfer $\hat a$ in basis $\Phimat$);
prior $R_\theta$ and heat symbols $\hat k_t$; weights $w_p, w_y, \sigma_y$; levels $N$
\State $\x \gets \Phimat^{-1}\!\big(\hat k_N \odot \Phimat\y\big)$
\Comment{initialize: blur $\y$ to the top level}
\State $\hat{\y}\gets \Phimat\y$
\For{$t = N, N\!-\!1, \dots, 1$}
  \State $\bmu \gets R_\theta(\x,\,t)$
  \Comment{learned reverse step (image domain) $\to$ prior mean $\bmu_{t-1}$}
  \State $\widehat{\y}_{t-1} \gets \hat k_{t-1}\odot\hat{\y}$
  \Comment{scale-matched measurement $\K_{t-1}\y$, Eq.~\eqref{eq:scalematch}}
  \State $w_y(t) \gets \dfrac{\alpha_y}{\sigma_y^2}\Big(1-\dfrac{t-1}{N}\Big)$
  \State $\displaystyle \hat{\x}\gets
     \frac{w_p\,(\Phimat\bmu)\;+\;w_y(t)\,\overline{\hat a}\odot\widehat{\y}_{t-1}}
          {w_p\;+\;w_y(t)\,|\hat a|^2}$
  \Comment{closed-form MAP, Eq.~\eqref{eq:wiener}}
  \State $\x \gets \mathrm{clip}\big(\Phimat^{-1}\hat{\x},\,-1,\,1\big)$
\EndFor
\State \Return $\x$
\end{algorithmic}
\end{algorithm}

Each iteration is one network evaluation plus two transforms; the data step is
$\mathcal{O}(d\log d)$ and exact. The method never differentiates the prior.

\section{Numerical results}
\label{sec:results}

\paragraph{Setup.}
Held-out CelebA-HQ-256, $n=16$ images, noise
$\sigma_y=0.05$. Priors: \emph{COLD-DIFF} and \emph{IHDM-211M}. Baseline:
Diffusion Posterior Sampling (DPS). Metrics:
PSNR$\uparrow$, SSIM$\uparrow$, LPIPS$\downarrow$, measurement consistency
$\mathrm{MC}=\|\y-\A\hat\x\|/\|\y\|\;\downarrow$.

\paragraph{Gaussian / heat deblur.}
With $\A$ a DCT-diagonal Gaussian blur, $\A$ and the heat prior
share the DCT basis, so~\eqref{eq:wiener} is exact. The non-hot HQS prior wins distortion
by a wide margin (Table~\ref{tab:gauss}); DPS wins perceptual LPIPS. Both non-hot priors
also exceed a best-$\lambda$ classical Wiener filter; Wiener's
deconvolution ringing gives it the worst LPIPS. Its
$\lambda$ is chosen by oracle test-set PSNR, i.e.\ Wiener's strongest setting.

\begin{table}[h]
\centering
\caption{Gaussian deblur ($\sigma_{\text{blur}}=4$), held-out CelebA-HQ-256, across noise
$\sigma_y$. Bold = best PSNR and best LPIPS per noise block.}
\label{tab:gauss}
\begin{tabular}{llcccc}
\toprule
$\sigma_y$ & Method & PSNR & SSIM & LPIPS & MC \\
\midrule
\multirow{6}{*}{0.05}
 & Observation & 22.85 & 0.529 & 0.565 & 0.122 \\
 & Wiener      & 25.44 & 0.698 & 0.585 & 0.094 \\
 & HQS + TV    & 26.04 & 0.742 & 0.364 & 0.093 \\
 & cold + HQS  & 25.80 & 0.692 & 0.442 & 0.092 \\
 & IHDM + HQS  & \textbf{27.13} & \textbf{0.778} & 0.360 & 0.093 \\
 & DPS         & 23.84 & 0.678 & \textbf{0.175} & 0.102 \\
\midrule
\multirow{6}{*}{0.10}
 & Observation & 21.48 & 0.352 & 0.719 & 0.198 \\
 & Wiener      & 24.63 & 0.671 & 0.645 & 0.185 \\
 & HQS + TV    & 25.23 & 0.714 & 0.408 & 0.183 \\
 & cold + HQS  & 25.11 & 0.683 & 0.460 & 0.181 \\
 & IHDM + HQS  & \textbf{26.47} & \textbf{0.756} & 0.380 & 0.182 \\
 & DPS         & 23.28 & 0.658 & \textbf{0.186} & 0.191 \\
\midrule
\multirow{6}{*}{0.20}
 & Observation & 18.53 & 0.188 & 0.927 & 0.349 \\
 & Wiener      & 23.41 & 0.613 & 0.747 & 0.344 \\
 & HQS + TV    & 23.44 & 0.662 & 0.481 & 0.348 \\
 & cold + HQS  & 23.99 & 0.664 & 0.474 & 0.341 \\
 & IHDM + HQS  & \textbf{25.41} & \textbf{0.724} & 0.396 & 0.342 \\
 & DPS         & 22.57 & 0.632 & \textbf{0.201} & 0.351 \\
\bottomrule
\end{tabular}
\end{table}

\paragraph{Motion deblur (realistic reflect boundaries, spatial CG).}
A Levin09 kernel~\#1 is \emph{asymmetric}, so it is not diagonalized by the DCT; and on a finite
image with a realistic \textbf{reflect (symmetric) boundary}---the standard blur-synthesis
convention (reflect-pad by half the kernel $\to$ convolve $\to$ crop)---the blur operator is block
Toeplitz-plus-Hankel, \emph{not} circulant, so no single transform diagonalizes it. A closed-form
\emph{frequency-domain} data step (the Wiener update~\eqref{eq:wiener}) would apply an inverse DFT
over the whole image and thereby re-impose periodic boundaries, producing border ringing on
realistic data. We therefore keep the same per-step MAP objective~\eqref{eq:map} but solve its data
step in the \emph{spatial} domain by conjugate gradient. Writing $\y=\mathbf H\x_0+\bm n$ with
$\mathbf H$ the reflect-boundary convolution matrix, each step solves the normal equations
\begin{equation}
\big(w_p\I + w_y\,\mathbf H^{\!\top}\mathbf H\big)\,\x \;=\; w_p\bmu + w_y\,\mathbf H^{\!\top}(\K_{t-1}\y)
\label{eq:cg}
\end{equation}
by CG ($\sim\!12$ iterations, matrix-free: each iteration is one $\mathbf H$ and one
$\mathbf H^{\!\top}$ spatial convolution; the adjoint $\mathbf H^{\!\top}$ is taken exactly by
autograd, verified $\langle\mathbf H\x,\mathbf r\rangle=\langle\x,\mathbf H^{\!\top}\mathbf r\rangle$
to $10^{-15}$). No iFFT over the image is formed, so there is no periodic wrap and no ringing, and
the cost is negligible next to the network evaluation. Table~\ref{tab:motion} reports every method
on identical reflect-boundary observations: the learned priors with the spatial-CG data step
(cold$+$CG, IHDM$+$CG) recover the full frame cleanly, whereas the frequency-domain baselines
(oracle Wiener, HQS$+$TV) and DPS inherit the boundary ring on realistic data. IHDM$+$CG is the
best full-frame method at every noise level ($29.74/28.40/26.78$\,dB PSNR), $3$--$5$\,dB above every
baseline: the circulant Wiener and HQS$+$TV data steps ring on the non-periodic data (HQS$+$TV reaches
only $22.66/24.82/24.54$\,dB) and DPS trails on distortion ($24.90/24.49/23.74$\,dB), though DPS wins
perceptual LPIPS at the two higher noise levels. The weak cold prior stays weak ($\approx\!22$\,dB),
as it was under HQS.

\begin{table}[h]
\centering
\caption{Motion deblur (Levin09~\#1) on realistic \textbf{reflect-boundary} observations, full
frame ($256^2$), across noise $\sigma_y$, held-out CelebA-HQ-256, $n=16$. Learned priors use the
spatial-CG data step (Eq.~\ref{eq:cg}); Wiener, HQS$+$TV and DPS are frequency-domain / circulant
baselines that inherit the boundary ring. Bold = best PSNR and best LPIPS per noise block.}
\label{tab:motion}
\begin{tabular}{llcccc}
\toprule
$\sigma_y$ & Method & PSNR & SSIM & LPIPS & MC \\
\midrule
\multirow{6}{*}{0.05}
 & Observation & 24.25 & 0.581 & 0.358 & 0.117 \\
 & Wiener      & 23.83 & 0.581 & 0.332 & 0.125 \\
 & HQS + TV    & 22.66 & 0.738 & 0.225 & 0.126 \\
 & cold + CG   & 22.66 & 0.440 & 0.491 & 0.061 \\
 & IHDM + CG   & \textbf{29.74} & 0.808 & \textbf{0.119} & 0.081 \\
 & DPS         & 24.90 & 0.718 & 0.168 & 0.122 \\
\midrule
\multirow{6}{*}{0.10}
 & Observation & 22.49 & 0.396 & 0.536 & 0.193 \\
 & Wiener      & 22.07 & 0.432 & 0.499 & 0.194 \\
 & HQS + TV    & 24.82 & 0.726 & 0.258 & 0.191 \\
 & cold + CG   & 22.05 & 0.399 & 0.542 & 0.142 \\
 & IHDM + CG   & \textbf{28.40} & 0.795 & 0.246 & 0.172 \\
 & DPS         & 24.49 & 0.701 & \textbf{0.173} & 0.198 \\
\midrule
\multirow{6}{*}{0.20}
 & Observation & 19.07 & 0.219 & 0.790 & 0.342 \\
 & Wiener      & 19.79 & 0.333 & 0.681 & 0.352 \\
 & HQS + TV    & 24.54 & 0.692 & 0.333 & 0.338 \\
 & cold + CG   & 22.10 & 0.419 & 0.503 & 0.302 \\
 & IHDM + CG   & \textbf{26.78} & 0.761 & 0.327 & 0.330 \\
 & DPS         & 23.74 & 0.673 & \textbf{0.184} & 0.347 \\
\bottomrule
\end{tabular}
\end{table}

\begin{figure}[htbp]
\centering
\includegraphics[width=\linewidth]{../results/four_way/figure_gaussian.png}
\caption{Gaussian deblur, held-out CelebA-HQ-256 (per-column mean PSNR/SSIM/LPIPS in the
headers). With $\A$ and the heat prior sharing the DCT basis, the non-hot HQS priors (cold,
IHDM) win distortion; DPS wins perceptual LPIPS.}
\label{fig:gauss}
\end{figure}

\begin{figure}[htbp]
\centering
\includegraphics[width=\linewidth]{../results/motion/figure_dct_vs_dft.png}
\caption{Motion deblur (Levin09~\#1) at $\sigma_y=0.05$ on realistic \textbf{reflect-boundary}
observations, full frame: HQS$+$TV, cold$+$CG and IHDM$+$CG plus DPS. The frequency-domain
baselines (oracle Wiener, HQS$+$TV) and DPS show the periodic \emph{banding} of a circulant
data step on non-periodic data, whereas the spatial-CG learned prior (IHDM$+$CG) reconstructs
the full frame cleanly---no iFFT, no wrap (Table~\ref{tab:motion}).}
\label{fig:dft}
\end{figure}

\paragraph{Noise robustness.}
We repeat the full comparison at higher noise, $\sigma_y\in\{0.10,0.20\}$
(Tables~\ref{tab:gauss},~\ref{tab:motion}, Figs.~\ref{fig:gauss_noise},~\ref{fig:motion_noise_full};
TV and Wiener re-tuned per noise). The rankings are stable. On Gaussian, IHDM$+$HQS keeps the best
PSNR/SSIM at every level ($27.13\!\to\!26.47\!\to\!25.41$\,dB) and DPS the best LPIPS. On motion
(reflect boundaries), IHDM$+$CG is the best full-frame method at every noise level, while DPS
increasingly wins LPIPS as noise grows. All methods lose $\sim\!2$--$4$\,dB from $\sigma_y{=}0.05$
to $0.20$, but the ordering is preserved --- the structural conclusions are not artifacts of a
single noise level.

\begin{figure}[htbp]\centering
\includegraphics[width=\linewidth]{../results/four_way_n0p10/figure.png}\\[2pt]
\includegraphics[width=\linewidth]{../results/four_way_n0p20/figure.png}
\caption{Gaussian deblur at $\sigma_y=0.10$ (top) and $0.20$ (bottom).}
\label{fig:gauss_noise}
\end{figure}

\begin{figure}[htbp]\centering
\includegraphics[width=\linewidth]{../results/motion_n0p10/figure_full.png}\\[2pt]
\includegraphics[width=\linewidth]{../results/motion_n0p20/figure_full.png}
\caption{Motion deblur, full frame, at $\sigma_y=0.10$ (top) and $0.20$ (bottom).}
\label{fig:motion_noise_full}
\end{figure}

\paragraph{Summary.}
For \textbf{Gaussian} blur, the measurement operator and the heat prior share the DCT
eigenbasis, so the non-hot HQS prior gives an exact, backprop-free, closed-form data step and is
the best distortion method---beating both a well-tuned classical baseline (HQS$+$TV) and DPS
(Table~\ref{tab:gauss}). For \textbf{motion} blur under realistic reflect boundaries the operator
is not diagonalized by any fast transform, so a closed-form frequency-domain (iFFT) data step
re-imposes periodic boundaries and rings; we instead solve the same per-step MAP objective with a
spatial \textbf{conjugate-gradient} data step (the true reflect operator $\mathbf H$,
Eq.~\ref{eq:cg}), which eliminates the border ring and makes IHDM$+$CG the best full-frame method
at every noise level (Table~\ref{tab:motion}). Two caveats keep the
picture honest: (i) \textbf{HQS$+$TV is a very strong classical prior} on Gaussian blur: within
the identical HQS loop a TV regularizer comes within $\sim\!1$\,dB of the learned prior on the
Gaussian full frame; on realistic reflect-boundary motion, however, its frequency-domain data
step rings like the other circulant solvers, and the spatial-CG learned prior wins; and
(ii) \textbf{DPS} is the most
robust choice when the operator's structure is unknown or perceptual LPIPS is the target,
since it never diagonalizes $\A$ and only needs $\A,\A^\top$. The learned diffusion prior
wins where its structural assumptions hold; TV and DPS are the baselines to beat elsewhere.

\end{document}
